\chapter{🎁 授受表現 — Dar y Recibir}
\unidadheader{23}{授受表現}{Dar y Recibir}{Expresar intercambio de objetos y favores}

\begin{tcolorbox}[vocabbox, title={📌 Vocabulario Nuevo}]
\begin{tabularx}{\linewidth}{llX}
\toprule
\textbf{Japonés} & \textbf{Español} & \textbf{Tipo} \\
\midrule
\vocabitem{プレゼント}{—}{regalo}{sust.}
\vocabitem{花}{はな}{flor}{sust.}
\vocabitem{手紙}{てがみ}{carta}{sust.}
\vocabitem{お土産}{おみやげ}{souvenir / regalo de viaje}{sust.}
\vocabitem{贈る}{おくる}{enviar como regalo}{v. gr.1}
\vocabitem{受け取る}{うけとる}{recibir / aceptar}{v. gr.1}
\vocabitem{喜ぶ}{よろこぶ}{alegrarse}{v. gr.1}
\vocabitem{感謝する}{かんしゃする}{agradecer}{v. irr.}
\vocabitem{お世話になる}{おせわになる}{recibir cuidado/ayuda de}{expr.}
\vocabitem{返す}{かえす}{devolver}{v. gr.1}
\bottomrule
\end{tabularx}
\end{tcolorbox}

\begin{tcolorbox}[phrasebox, title={🔄 Diagrama: あげる / もらう / くれる}]
\begin{verbatim}
  [Yo] ---あげます---> [Otro]   (yo doy a otro)
  [Yo] <--もらいます--- [Otro]   (yo recibo de otro)
  [Yo] <---くれます---- [Otro]   (otro da a mí)

  Clave: くれる siempre tiene a 私 (yo/mi grupo) como receptor.
\end{verbatim}
\end{tcolorbox}

\subsection*{🗣️ Frases Clave}

\frase{\ruby{友}{とも}だちにプレゼントをあげました。}
  {Le di un regalo a mi amigo.}
\frase{お\ruby{母}{かあ}さんからセーターをもらいました。}
  {Recibí un suéter de mi mamá.}
\frase{\ruby{田中}{たなか}さんが\ruby{花}{はな}をくれました。}
  {El Sr. Tanaka me dio flores.}
\frase{〜してもらえますか？}{¿Me podrías hacer el favor de ~?}

\subsection*{⚙️ Gramática}

\patron{Dar a alguien: あげます}
  {A が B に もの を あげます}
  {\ej{\ruby{友達}{ともだち}に\ruby{本}{ほん}をあげました。}
    {Le di un libro a mi amigo.}}

\patron{Recibir de alguien: もらいます}
  {A が B に/から もの を もらいます}
  {\ej{せんせいに/からアドバイスをもらいました。}
    {Recibí un consejo del profesor.}}

\patron{Alguien me da: くれます}
  {B が A に もの を くれます (A = yo o mi grupo)}
  {\ej{\ruby{彼女}{かのじょ}がチョコをくれました。}
    {Mi novia me dio chocolate.}}

\patron{Favores con て形}
  {V-て + あげる / もらう / くれる}
  {\ej{\ruby{宿題}{しゅくだい}を\ruby{手伝}{てつだ}ってあげました。}
    {Le ayudé con su tarea.}
  \ej{\ruby{写真}{しゃしん}を\ruby{撮}{と}ってもらえますか？}
    {¿Me podrías tomar una foto?}}

\begin{tcolorbox}[ejerciciobox, title={📝 Ejercicios Rápidos}]
\begin{enumerate}
  \item ¿あげる o くれる? Tu amigo te da un libro → ?
  \item Pide un favor usando てもらえますか.
  \item Traduce: \ruby{母}{はは}が\ruby{料理}{りょうり}を\ruby{作}{つく}ってくれました。
\end{enumerate}
\vspace{0.4em}
\textbf{Respuestas:}
1) くれる (\ruby{友達}{ともだち}が\ruby{本}{ほん}をくれました。)\quad
2) (personal)\quad
3) Mi mamá me cocinó (me hizo el favor de cocinar).
\end{tcolorbox}

\begin{tcolorbox}[kanjibox, title={🀄 Kanjis de la Unidad}]
\begin{tabularx}{\linewidth}{cllXX}
\toprule
\textbf{Kanji} & \textbf{On} & \textbf{Kun} & \textbf{Significado} & \textbf{Vocab} \\
\midrule
\kanjirow{花}{か}{はな}{flor}{\ruby{花}{はな} / \ruby{花火}{はなび}}
\kanjirow{手}{しゅ}{て}{mano}{\ruby{手紙}{てがみ} / \ruby{手伝う}{てつだう}}
\kanjirow{紙}{し}{かみ}{papel}{\ruby{手紙}{てがみ} / \ruby{紙}{かみ}}
\kanjirow{贈}{ぞう}{おく}{regalar}{\ruby{贈}{おく}る / \ruby{贈り物}{おくりもの}}
\kanjirow{受}{じゅ}{う}{recibir}{\ruby{受}{う}ける / \ruby{受け取る}{うけとる}}
\bottomrule
\end{tabularx}
\end{tcolorbox}
